\selectlanguage{ngerman}%
\frenchspacing%
\begin{ledgroupsized}[r]{120mm}%
\footnotesize%
\pstart%
\noindent%
\textbf{Überlieferung:}%
\pend%
\end{ledgroupsized}%
\begin{ledgroupsized}[r]{114mm}%
\footnotesize%
\pstart%
\parindent -6mm%
\makebox[6mm][l]{\textit{L}}%
Auszüge mit Bemerkungen aus
\textsc{D.\;Papin}, \glqq Nova methodus ad vires motrices validissimas levi pretio comparandas\grqq\ (\textit{AE}, August 1690, S.\;410--414):\cite{01023}\cite{01479}
LH XXXVIII Bl.\;156.
Ein Blatt 8\textsuperscript{o};
Ränder zum Teil beschnitten;
\href{\Scan{LH_38_156wz}}{Fragment} eines Wasserzeichens.
Zwei Seiten;
Textfolge (nicht von Leibniz festgelegt):
Bl.\;\href{\Scan{LH_38_156v}}{156\,v\textsuperscript{o}},
Bl.\;\href{\Scan{LH_38_156r}}{156\,r\textsuperscript{o}}.
\pend%
\end{ledgroupsized}%
\vspace{5mm}%
\begin{ledgroup}%
\footnotesize%
\pstart%
\noindent%
\textbf{Datierungsgründe:}
Die Datierung der Auszüge N.\;\ref{55804} beruht auf %
inhaltlichen Betrachtungen.
Siehe die ausführliche Begründung auf S.\;\pageref{55804_dat}.%
\pend%
\end{ledgroup}%
\selectlanguage{latin}%
\frenchspacing%
\vspace*{8mm}
\normalsize%
\pstart%
\noindent%
\lbrack%
\href{\Scan{LH_38_156v}}%
{156\,v\textsuperscript{o}}%
\rbrack\ %
\edtext{}{%
\xxref{LH_38_156_156v_hkm-1}{LH_38_156_156v_hkm-2}%
\lemma{Si pulverem \lbrack...\rbrack\ diminuatur}%
\Cfootnote{%
Vgl.
\textsc{D.\;Papin},
\glqq Nova methodus ad vires motrices validissimas levi pretio comparandas\grqq\
(\textit{AE}, August 1690, S.\;410).\cite{01023}\cite{01479}%
}}%
\edtext{%
\edlabel{LH_38_156_156v_hkm-1}%
Si pulverem pyrium%
\indexS{pulvis pyrius accensus in tubo}
in Tubo accendas,
oritur in eo%
}{%
\lemma{\lbrack\href{\Scan{LH_38_156v}}{156\,v\textsuperscript{o}}\rbrack}%
\Bfootnote{%
\textit{(1)}~\textit{Pulvis pyrius in infima Tubi parte accensus totam cavitatem} Tubi it%
\indexS{cavitas tubi} %
\textit{(2)}~Si pulverem %
\lbrack...\rbrack\ in eo%
~\textit{L}%
}}
vis quaedam Elastica%
\indexS{vis elastica e pulvere pyrio}
\edtext{%
quae Tubum implet%
\indexS{tubus impletus a vi elastica}
aeremque inclusum%
\indexS{aer inclusus in tubo}
expellit%
\indexS{aer expulsus a tubo} %
atque eadem opera embolum%
\indexS{embolus tubi}
attollit.
Sed eadem vis %
extincta flamma%
\indexS{flamma extincta}%
}{%
\lemma{quae}%
\Bfootnote{%
\textit{(1)}~embolum attollit, %
\textit{(2)}~Tubum implet %
\textit{(a)}~et embolum in %
\textit{(b)}~aeremque inclusum %
\lbrack...\rbrack\ embolum attollit. %
\textit{(aa)}~Qui deinde %
\textit{(bb)}~Sed eadem  vis extincta flamma%
~\textit{L}%
}}
subito iterum perit,%
\indexS{vis periens}
tubusque vacuus relinquitur,%
\indexS{tubus vacuus aere}
et pondus aeris externi%
\indexS{pondus aeris externi}%
\indexS{aer externus}
rursus embolum%
\indexS{embolus tubi}
\edtext{%
intrudit,
simulque%
}{%
\lemma{intrudit}%
\Bfootnote{%
\textit{(1)}~eademque opera %
\textit{(2)}~simulque%
~\textit{L}%
}}
pondus embolo annexum%
\indexS{pondus annexum embolo}
satis magnum attollit.
Interim fatendum
\edtext{%
est
non posse perfectum aeris vacuum%
\indexS{vacuum aeris perfectum}%
\indexS{aer evacuatus}
hac arte obtineri,%
\indexS{vacuum obtentum arte}%
}{%
\lemma{est}%
\Bfootnote{%
\textit{(1)}~duo incommoda vim nonnihil minuere %
\textit{(2)}~non posse %
\lbrack...\rbrack\ arte obtineri,%
~\textit{L}%
}}
sed quintam circiter aeris partem intus restare.
Unde dimidius tantum obtinetur effectus%
\indexS{effectus dimidius} %
\edtext{%
et pro pondere 300 librarum%
\indexS{libra} %
tantum attolli possunt 150
ad pedis altitudinem;%
\indexS{altitudo elevationis}%
\indexS{pes}%
}{%
\lemma{et pro}%
\Bfootnote{%
\lbrack...\rbrack\ pedis altitudinem %
\textit{erg.\;L}%
}}
et praeterea inaequalis est actio Machinae,%
\indexS{actio machinae}%
\indexS{machina}
quia aer ille
qui intus superfuerat
descendente embolo%
\indexS{embolus tubi}
magis magisque
ad densitatem naturalem%
\indexS{densitas aeris naturalis}
\edtext{%
redit,
atque%
}{%
\lemma{redit,}%
\Bfootnote{%
\textit{(1)}~et %
\textit{(2)}~atque%
~\textit{L}%
}}
adeo continue diminuitur pressio aeris externi.%
\indexS{pressio aeris externi}
Unde necesse est
machinam ita applicatam esse oneri,%
\indexS{machina applicata oneri}%
\indexS{onus}
ut resistentia itidem continue%
\indexS{resistentia in machina}
\edlabel{lh038_156v_1a}%
\edlabel{LH_38_156_156v_hkm-2}%
diminuatur.%
\edtext{}{%
\xxref{lh038_156v_1a}{lh038_156v_1b}%
\lemma{diminuatur.}%
\Bfootnote{%
\textit{(1)}~Alia %
\textit{(2)}~Loco v %
\textit{(3)}~Caret his incommodis alia%
~\textit{L}%
}}
\pend%
\newpage%
\pstart%
\edtext{}{%
\xxref{LH_38_156_156v_rto-1}{LH_38_156_156v_rto-2}%
\lemma{Caret \lbrack...\rbrack\ respondens}%
\Cfootnote{%
a.a.O. (S.\;411).\cite{01023}\cite{01479}%
}}%
\edlabel{LH_38_156_156v_rto-1}%
Caret his incommodis%
\indexS{incommodum}
alia%
\edlabel{lh038_156v_1b}
quaedam Vis Elastica%
\indexS{vis elastica e vapore}
quam exercet aqua in vaporem%
\indexS{vapor aqueus}%
\indexS{aqua resoluta in vapore}
calore%
\indexS{calor resolvens aquam in vaporem}
\edtext{%
resoluta.
Hic enim vapor%
\indexS{vapor aqueus}
itidem vim habet%
\indexS{vis vaporis}
embolum expellendi;%
\indexS{vis expellendi embolum}%
\indexS{embolus tubi}
et quod caput rei est,
refrigerato tubo%
\indexS{tubus refrigeratus}%
\indexS{refrigeratio tubi}%
}{%
\lemma{resoluta}%
\Bfootnote{%
\textit{(1)}~; quae et ipsa superveniente mox frigore  %
\textit{(2)}~; qui %
\textit{(3)}~. Hic enim %
\lbrack...\rbrack\ embolum expellendi; %
\textit{(a)}~Tubumque ite %
\textit{(b)}~et quod %
\lbrack...\rbrack\ refrigerato tubo%
~\textit{L}%
}}
in aquam redit,%
\indexS{vapor reditus in aquam}%
\indexS{aqua e vapore}
nullo aeris vestigio relicto,%
\indexS{vestigium aeris relictum}
atque ita perfectum efficitur vacuum aeris%
\indexS{vacuum aeris perfectum} %
\edtext{%
visque pressioni atmosphaerae aequabilis%
\indexS{vis aequabilis pressioni atmosphaerae}%
\indexS{pressio atmosphaerae}%
\indexS{atmosphaera} %
et magnitudini Tubi respondens.%
\indexS{vis respondens magnitudini tubi}%
\indexS{tubus vacuus aere}%
\edlabel{LH_38_156_156v_rto-2}%
}{%
\lemma{visque}%
\Bfootnote{%
\textit{(1)}~pressione atmosphaerae respondens %
\textit{(2)}~pressioni atmosphaerae %
\lbrack...\rbrack\ Tubi respondens.%
~\textit{L}%
}}
(+\;%
Adde quod Tubus hoc modo semper manet clausus,
neque ut in accensione pulveris pyrii%
\indexS{accensio pulveris pyrii}%
\indexS{pulvis pyrius}
alternatim pro novi pulveris ingressu%
\indexS{ingressus pulveris pyrii in tubum} %
aperiri postulat.
Sumtus etiam minores sunt,
et parabilis atque
\edtext{%
ubique semper obvia%
}{%
\lemma{ubique}%
\Bfootnote{%
\textit{(1)}~obvia %
\textit{(2)}~semper obvia%
~\textit{L}%
}}
tantum materia desideratur,%
\indexS{materia parabilis}%
\indexS{materia obvia}
Aqua scilicet et
\edtext{%
ignis,%
\indexS{ignis}
sed promtitudo operationis%
\indexS{promtitudo operationis}%
\indexS{operatio machinae}%
}{%
\lemma{ignis,}%
\Bfootnote{%
\textit{(1)}~illud l %
\textit{(2)}~sed promtitudo operationis%
~\textit{L}%
}}
utique minor est.%
\;+)
\lbrack%
\href{\Scan{LH_38_156r}}%
{156\,r\textsuperscript{o}}%
\rbrack%
\pend%
\pstart%
\edtext{%
Experimentum%
\indexS{experimentum de vi vaporis}
sumtum est in Tubo metallico tenui%
\indexS{tubus metallicus tenuis}
\textit{%
cujus diameter $2\frac{1}{2}$ digitos non excedit.%
\indexS{digitus}%
}
Is 60 libras%
\indexS{libra}
attollere potest aequabili velocitate.%
\indexS{velocitas aequabilis}
Pondus ipsius vix est quinque unciarum%
\indexS{uncia}
(+\;%
altitudo, ni fallor, 4 pedum%
\;+).
Puto tubum
qui
\edtext{%
vix%
}{%
\lemma{vix}%
\Bfootnote{%
\textit{erg.\;L}%
}}
40 libras aequet,
bis mille libras singulis operationibus
ad 4 pedum altitudinem elevare posse.%
\indexS{altitudo elevationis}
Porro tubus ille exiguus
in quo factum est experimentum%
\indexS{experimentum de vi vaporis}%
\lbrack,\rbrack\
mediocri igne%
\indexS{ignis}
ita calefactus est,%
\indexS{tubus calefactus}
ut tempore unius minuti%
\indexS{minutum}
Embolus%
\indexS{embolus tubi}
ad summum attolleretur.
Nec majori tempore opus est%
\indexS{tempus elevationis emboli}
pro majoribus tubis,%
\indexS{tubus calefactus}
modo ignis%
\indexS{ignis}
proportione respondeat.
Notum est
columnam aeris incumbentem%
\indexS{columna aeris}%
\indexS{aer incumbens}
\edtext{%
Tubo%
\indexS{tubus calefactus}
unius pedis%
\indexS{pes}%
}{%
\lemma{Tubo}%
\Bfootnote{%
\textit{(1)}~40 pedum %
\textit{(2)}~unius pedis%
~\textit{L}%
}}
diametrum habenti
aequare bis mille libras.%
\indexS{libra}
Itaque si diameter sit 2 pedum,%
\indexS{pes}
poterunt tempore unius minuti%
\indexS{tempus elevationis emboli}%
\indexS{minutum}
attolli 8000 librae ad
\edtext{%
altitudinem 4 pedum%
}{%
\lemma{altitudinem}%
\Bfootnote{%
\textit{(1)}~40 pe %
\textit{(2)}~pe %
\textit{(3)}~4~pedum%
~\textit{L}%
}}
(+\;seu 800 librae
ad altitudinem 40 pedum\;%
\lbrack+)\rbrack. 
\edtext{%
Adeoque non minor est effectus%
\indexS{effectus}
quam si 80 librae attollentur%
\indexS{libra}%
}{%
\lemma{Adeoque}%
\Bfootnote{%
\textit{(1)}~8 secundis %
\textit{(a)}~atto %
\textit{(b)}~librae %
\textit{(2)}~non minor %
\lbrack...\rbrack\ librae attollentur%
~\textit{L}%
}}
ad altitudinem%
\indexS{altitudo elevationis emboli}
\edtext{%
pedum~40;
seu $106\frac{2}{3}$%
}{%
\lemma{pedum 40;}%
\Bfootnote{%
\textit{(1)}~($80 \cdot 40=30 \cdot \pi $, seu $\pi = 320 : 3 = 106 \frac{2}{3}$) %
\textit{(2)}~seu $106 \frac{2}{3}$%
~\textit{L}%
}}
libras,%
\indexS{libra}
vel centenarium,
ad altitudinem%
\indexS{altitudo elevationis emboli}
pedum 30,%
\indexS{pes}
\edtext{%
adeoque sexagies%
}{%
\lemma{adeoque}%
\Bfootnote{%
\textit{(1)}~sexcenties %
\textit{(2)}~sexagies%
~\textit{L}%
}}
in hora una%
\indexS{hora}
\edtext{%
repeti posset%
}{%
\lemma{repeti}%
\Bfootnote{%
\textit{(1)}~poterit %
\textit{(2)}~posset%
~\textit{L}%
}}
operatio.%
\indexS{operatio machinae}
Sufficit autem
infundi tubo%
\indexS{tubus calefactus}
exiguam aquae quantitatem%
\indexS{aqua infusa tubo}
ad 3 vel 4 linearum altitudinem.
Ignis continuus,%
\indexS{ignis continuus}
ne
\edtext{%
frustra consumatur%
}{%
\lemma{frustra}%
\Bfootnote{%
\textit{(1)}~agat %
\textit{(2)}~consumatur%
~\textit{L}%
}}
durante emboli intrusione%
\indexS{intrusio emboli}%
\indexS{embolus intrusus}%
\lbrack,\rbrack\
de tubo in tubum%
\indexS{tubus calefactus}
circumferri potest in fornace%
\indexS{fornax}
ex tenui ferro confecta.%
\indexS{ferrum tenue}%
}{%
\lemma{Experimentum \lbrack...\rbrack\ confecta}%
\Cfootnote{%
a.a.O. (S.\;411\,f.).\cite{01023}\cite{01479}
Leibniz hat hier stark umformuliert und eigenständig umgerechnet.%
}}
\pend%
\newpage%
\pstart%
\edtext{%
\textit{%
Vis haec applicari potest%
\indexS{vis vaporis}%
\indexS{vis applicata}
ad extrahendam ex fodinis aquam%
\indexS{aqua extracta e fodinis}%
\indexS{fodina}
aut mineram,%
\indexS{minera extracta e fodinis}
ferreos globos ad maximam distantiam projiciendos,%
\indexS{globus ferreus projiciendus}
naves adverso vento%
\indexS{ventus adversus}%
}
vel nullo
\textit{%
provehendas%
\indexS{navis provehenda}%
}
remigum loco.%
\indexS{remex}
Neque enim ut remiges gravat%
\indexS{remex}
(+\;%
imo magis ob alimentum ignis%
\indexS{alimentum ignis}%
\indexS{ignis continuus}
continuum%
\;+)
nec tantum loci occupat %
(+\;%
eadem difficultas%
\;+)%
\indexS{difficultas}
minoreque sumtu indiget%
\lbrack;\rbrack\
possent adhiberi remi rotatiles%
\indexS{remus rotatilis}
\edtext{%
exemplo triremis%
\indexS{triremis}
principis Roberti%
\indexN{\textls{Pfalz}, Ruprecht von der (Robertus Palatinus, princeps Robertus) 1619\textendash1682}%
}{%
\lemma{exemplo}%
\Bfootnote{%
\textit{(1)}~principis Roberti %
\textit{(2)}~triremis principis Roberti%
~\textit{L}%
}%
\lemma{exemplo \lbrack...\rbrack\ Roberti}%
\Cfootnote{%
Ruprecht von der Pfalz war 1672--1679 erster Admiral der englischen Marine.
Papin erwähnt, dessen von Pferdekraft betriebenes Ruderschiff in London\indexO{London} gesehen zu haben (a.a.O, S.\;413).\cite{01479}%
}%
}
quam equi%
\indexS{equus}
intus
\edtext{%
agitabant.
Haec%
}{%
\lemma{agitabant.}%
\Bfootnote{%
\textit{(1)}~Quae %
\textit{(2)}~Haec%
~\textit{L}%
}}
cymbam regiam%
\indexS{cymba}
sedecim remigum%
\indexS{remex}
longe post se relinquebat%
}{%
\lemma{\textit{Vis} \lbrack...\rbrack\ relinquebat}%
\Cfootnote{%
a.a.O. (S.\;412\,f.).\cite{01023}\cite{01479}%
}}
(+\;%
ad currus agendos,%
\indexS{currus}
\edtext{%
ad \bracks{navigia}%
\indexS{navigium}
\pgrk{Peri'androu}%
}{%
\lemma{ad}%
\Bfootnote{%
\bars{navigia \textit{gestr., wieder gültig macht Hrsg.}} \pgrk{Peri'androu}%
~\textit{L}%
}}
\edtext{%
\pgrk{Kle'anjhs}%
\;+).
}{%
\lemma{\pgrk{Peri'androu Kle'anjhs}}%
\Cfootnote{%
Anspielung unklar.
}}
\pend%
\pstart%
(+\;Majora fortasse praestari possunt
aere compresso calefacto.%
\indexS{aer compressus calefactus}%
\indexS{aer calefactus}%
\;\lbrack+)\rbrack%
\pend%
\vspace{4.0em}%
\pstart%
\noindent%
\centering%
*     *     *
\pend%
\vspace{3.0em}%
\selectlanguage{german}%
\frenchspacing%
\footnotesize%
\pstart%
\noindent%
\textls{Ausführliche Datierungsgründe:}%
\pagelabel{55804_dat}
Der undatierte und titellose Text N.\;\ref{55804} umfasst Leibnizens eigenhändige, stark zusammengefasste und mit eigenständigen Bemerkungen versehene Auszüge aus dem 1690 im Augustheft der \textit{Acta eruditorum}\cite{01023} veröffentlichten Aufsatz \glqq Nova methodus ad vires motrices validissimas levi pretio comparandas\grqq\ von Denis Papin.\cite{01479}\indexN{\textls{Papin} (Papinus), Denis 1647\textendash1713}
Unter der Annahme, dass Leibniz das Augustheft der \textit{Acta eruditorum}\cite{01023} kurz nach dessen Erscheinen hat lesen können, ist als Terminus post quem der Datierung August 1690 anzusehen.
\pend%
\pstart%
In seinem Brief an Papin vom 14.\;(24.) April 1698 spielt Leibniz eindeutig und zutreffend auf den besagten Aufsatz des französischen Naturforschers an (\textit{LSB} III,\;7 N.\;187, S.\;753.13--16).\cite{01480}
Demgemäß ist anzunehmen, dass er spätestens zu jenem Zeitpunkt Papins\indexN{\textls{Papin} (Papinus), Denis 1647\textendash1713} Aufsatz gelesen und exzerpiert hatte.
Mitte April 1698 ist folglich als Terminus ante quem der Datierung von N.\;\ref{55804} zu betrachten.
Das im Textträger vorliegende Fragment eines Wasserzeichens ist bei Redaktionsschluss nicht ermittelt worden.
\pend%
\selectlanguage{latin}%
\frenchspacing%
\normalsize%
\newpage%
